\section{Experimental Reinterpretation}


本節では、前節までに定式化したログ参照モデルを用いて、既存の代表的な意識実験を再解釈する。

ここで行うのは、実験結果の否定ではない。

むしろ、既存実験が示していたものを、意識生成モデルではなく、時間スケール間のログ参照モデルとして読み直すことである。

意識研究における多くの混乱は、処理、行動、報告、意識を単一の時間軸上に並べてしまうことから生じている。

この見方では、早く現れたものが原因であり、遅く現れたものは結果であると解釈されやすい。

しかし、本稿の枠組みでは、システムは単一の時計で動いていない。

高速層は予測し、低速層は沈殿し、可塑性ゲートは両者の間に遅延と粘性を導入する。

したがって、実験で観測される時間差や報告不能性は、意識の欠如を意味するとは限らない。

それらは、処理がログ参照可能な状態へ到達するまでの、時間的・経路的条件を示している。

\subsection{Libet-Type Experiments}

Libet 型実験では、被験者が任意のタイミングで手首や指を動かし、その際に「動かそうと思った瞬間」を報告する。

同時に、脳活動が測定される。

よく知られているように、運動に先立って準備電位が立ち上がり、その後で被験者の意識的な意図報告が現れる。

従来、この結果はしばしば次のように解釈されてきた。

脳が先に決め、意識は後からそれを知る。

あるいは、自由意志は幻想である。

しかし、この解釈は、決定、脳活動、意識報告を同一の時間軸上に置くことを前提としている。

本稿の枠組みでは、この前提を採らない。

準備電位は、高速層における行動候補の形成として読むことができる。

それは、未来志向の予測的・運動的ダイナミクスが、ある行動アトラクタへ落ち込み始めたことを示す信号である。

一方で、「いま動かそうと思った」という報告は、その高速処理が低速層のログ参照構造へ接続された後に現れる。

このとき初めて、処理は「自分がそうしようとしている」という報告可能な形式を取る。

したがって、Libet 型実験で観測される時間差は、意識の遅れではなく、速度スケール間の変換時間である。

高速層の処理速度を $v_f$、低速層の沈殿速度を $v_s$ とすれば、両者の速度差は次のように表される。

\begin{equation}
\Delta v = v_f - v_s
\end{equation}

この速度差は、内部遅延 $\tau$ として現れる。

\begin{equation}
\tau \propto \frac{\Delta v}{\mu}
\end{equation}

ここで $\mu$ は、高速層と低速層の結合粘性である。

つまり、Libet 型実験が測定したものは、自由意志の有無ではなく、高速な行動準備が低速な報告可能ログへ変換されるまでの遅延である。

この見方では、「意識は遅れている」という表現も十分ではない。

意識は、そもそも高速層と同じ時計で動いていない。

意識は、高速処理がログ参照可能な状態へ変換された地点で開かれる。

そのため、準備電位が先行することは、意識が無力であることを示さない。

それは、意識が行動準備そのものではなく、行動準備が低速層へ沈殿し、参照可能になる通過点であることを示している。

\subsection{Why Free Will Is Not the Proper Question Here}

Libet 型実験が繰り返し自由意志の議論へ引き込まれる理由は、「決定」が一点で起きるという前提にある。

しかし、本稿の構造では、決定は単一の時点に局在しない。

高速層では、複数の行動候補が予測的に立ち上がる。

可塑性ゲートは、その一部が低速層へ沈殿しうるかを調整する。

低速層では、過去ログ、意味勾配、身体状態、文脈が、報告可能な形式を作る。

この全体を一つの「決定点」に圧縮すると、時間構造が失われる。

そのため、Libet 型実験に対して「意識が決めたのか、脳が決めたのか」と問うことは、構造的には不適切である。

適切な問いは、次のようなものである。

どの処理が高速層で先行したのか。

その処理はどの程度の遅延 $\tau$ を経て低速層へ接続されたのか。

どのログが意味勾配 $\nabla \Phi$ によって参照されたのか。

どの時点で報告可能な形式が成立したのか。

このように問いを置き換えると、Libet 型実験は自由意志を否定する実験ではなく、単一時間軸モデルの限界を可視化した実験になる。

\subsection{Split-Brain Experiments}

分離脳実験では、脳梁を切断した患者において、左右半球の情報共有が制限される。

この状況で、右半球に提示された刺激に対して適切な行動反応が起きる一方で、その刺激を言語的に報告できない場合がある。

この現象は、しばしば「二つの意識」や「二つの自己」の証拠として語られてきた。

しかし、本稿の枠組みでは、この読みは過剰である。

分離脳実験が示しているのは、意識が二重化することではない。

それは、処理、沈殿、報告の経路が非対称化しうるということである。

右半球で処理が起きていても、その処理が左半球優位の言語報告経路へ到達しなければ、被験者はそれを明示的に語れない。

このとき、処理が存在しないわけではない。

行動反応が起きていることからも、局所的な処理は明らかに走っている。

しかし、その処理は報告可能な低速ログへ統合されていない。

したがって、分離脳実験で失われているのは、処理そのものではなく、参照経路である。

この点は、Libet 型実験と対応している。

Libet 型実験では、処理と報告の間に時間的遅延が現れる。

分離脳実験では、処理と報告の間に経路的断絶が現れる。

どちらも、処理が明示的意識になるためには、低速層へのログ沈殿と参照可能性が必要であることを示している。

\subsection{Reporting Bottleneck and Log Accessibility}

分離脳実験において特に重要なのは、報告可能性と意識を同一視しないことである。

報告可能性は、意識状態の一部を外部化する経路である。

しかし、報告できないことは、処理が起きていないことを意味しない。

本稿の枠組みでは、明示的報告には少なくとも三つの条件が必要である。

第一に、処理が高速層で生じていること。

第二に、その処理が低速層へ沈殿し、ログ参照可能になっていること。

第三に、そのログが言語または行動の報告経路へ接続されていること。

分離脳では、この第三の条件が非対称化する。

そのため、ある処理は行動として表出するが、言語報告としては現れない。

この非対称性を「二つの意識」と読む必要はない。

より正確には、ログ参照経路と報告経路の一部が切断されているのである。

\subsection{Rationalization as Log Completion}

分離脳実験でよく知られる現象として、左半球が、実際には参照できていない行動の理由を後から説明することがある。

これはしばしば「作話」や「合理化」と呼ばれる。

本稿の立場では、この合理化は意図的な虚偽ではない。

それは、欠落したログを補完し、低速層の一貫性を保とうとする自然な振る舞いである。

低速層は、参照可能なログを用いて自己の連続性を維持する。

しかし、処理経路が切断されている場合、行動は起きているのに、その行動を生んだログが参照できない。

このとき、システムは欠落をそのまま欠落として保持するより、参照可能なログからもっとも整合的な説明を生成する。

この合理化は、意識が嘘をついていることを示さない。

むしろ、意識が参照可能なログから一貫性を構成していることを示している。

この点で、分離脳実験は自我が主体であるという見方を弱める。

自我はすべてを知っている中心ではない。

それは、参照可能なログから自己連続性を組み立てる変換点である。

\subsection{Experimental Correspondence to the Model}

Libet 型実験と分離脳実験は、一見すると異なる問題を扱っている。

前者は時間差を扱い、後者は半球間の経路分断を扱う。

しかし、本稿の枠組みでは、両者は同じ構造の異なる断面である。

Libet 型実験は、次の関係を可視化している。

\begin{equation}
\text{fast preparation}
\rightarrow
\text{delayed sedimentation}
\rightarrow
\text{reportable timestamp}
\end{equation}

分離脳実験は、次の関係を可視化している。

\begin{equation}
\text{localized processing}
\rightarrow
\text{blocked transmission}
\rightarrow
\text{asymmetric reportability}
\end{equation}

どちらの場合も、処理そのものと意識報告は同一ではない。

処理は、可塑性ゲート、遅延、沈殿、報告経路を通過して初めて明示的意識として現れる。

この構造を用いれば、既存実験は次のように再配置される。

\begin{itemize}
\item Libet 型実験：速度スケール差と内部遅延 $\tau$ の露出。
\item 分離脳実験：ログ参照経路と報告経路の非対称化。
\item 合理化：参照不能ログの欠落に対する低速層の整合化。
\item 意識報告：処理そのものではなく、沈殿したログの外部化。
\end{itemize}

この再配置によって、意識は因果の源泉でも、単なる後付け幻想でもなくなる。

意識は、処理がログ参照可能な形へ変換されたときに開かれる状態である。

\subsection{What These Experiments Do Not Show}

最後に、これらの実験が示していないことを明確にしておく。

Libet 型実験は、自由意志が存在しないことを示したわけではない。

それは、行動準備と意識報告が異なる時間スケールで生じることを示した。

分離脳実験は、人間の内部に二つの自己が存在することを示したわけではない。

それは、処理経路と報告経路が切断されると、ログ参照が非対称化することを示した。

合理化は、意識が常に虚偽を語ることを示したわけではない。

それは、参照できないログの欠落を、低速層が一貫性のある形で補完することを示した。

したがって、これらの実験は意識を消し去るものではない。

むしろ、意識を正しい位置に戻す。

それは、行動の起点でも、脳内に生成される対象でもなく、速度スケール間のログ参照が報告可能な形に開かれた状態である。

\subsection{Summary}

本節では、Libet 型実験と分離脳実験を、本稿のログ参照モデルから再解釈した。

両者に共通するのは、処理と意識報告の非同一性である。

処理は走る。

しかし、それが低速層へ沈殿し、ログとして参照可能になり、報告経路へ接続されなければ、明示的意識としては現れない。

したがって、既存実験が示していたのは、意識の無力さではない。

それは、意識が成立するための時間的・経路的条件である。

次節では、このログ参照が沈殿境界でどのように現象的な前景を持つのか、すなわちクオリアを境界現象として位置づける。
